\begin{theorem}
For every real \(\varepsilon>0\), there is an integer \(n_0\) such that for
every integer \(n\ge n_0\) there exists a simple graph \(G\) on
\(\mathrm{Fin}(n)\) satisfying
\[
G.\mathrm{CliqueFree}(3)
\]
and
\[
\noderef{IndependenceNumberLess}
\left(G,(1+\varepsilon)\sqrt{(n:\mathbb R)\log(n:\mathbb R)}\right).
\]
Equivalently, for all sufficiently large \(n\) there is an \(n\)-vertex
triangle-free graph with
\(\alpha(G)<(1+\varepsilon)\sqrt{n\log n}\).
\end{theorem}
\begin{proof}
Fix \(\varepsilon>0\).  Apply \(\noderef{TwoBitePositiveExistence}\) with this
\(\varepsilon\).  It gives a threshold \(n_0\) such that, for every integer
\(n\ge n_0\), there exist an integer \(m\), a real number \(p\), and a sample
\[
\omega\in \mathrm{TwoBiteSample}(n,m,p)
\]
with the conclusions recorded in that child node.  The present theorem does
not use the displayed definitions of \(m\) and \(p\); it only needs the final
shadow graph and the two final graph properties supplied for that same
\(\omega\).

Fix \(n\ge n_0\), and choose \(m,p,\omega\) from
\(\noderef{TwoBitePositiveExistence}\).  Define
\[
G=(\noderef{TwoBiteFinalGraph}(\omega)).2.2 .
\]
By the definition of \(\noderef{TwoBiteFinalGraph}\), this third component is
the simple shadow graph on the vertex type \(\mathrm{Fin}(n)\).  Thus \(G\) is
a simple graph on \(\mathrm{Fin}(n)\), exactly as required in the statement.

We first prove \(G.\mathrm{CliqueFree}(3)\).  The no-triangle conjunct supplied
for \(\omega\) says exactly that there do not exist vertices
\(u,v,w\in\mathrm{Fin}(n)\) such that
\[
u\ne v,\qquad u\ne w,\qquad v\ne w
\]
and all three shadow adjacencies
\[
G.\mathrm{Adj}(u,v),\qquad G.\mathrm{Adj}(u,w),\qquad G.\mathrm{Adj}(v,w)
\]
hold.  This is the shadow-graph conclusion proved in
\(\noderef{TwoBiteDeletionTriangleFree}\) and included in
\(\noderef{TwoBitePositiveExistence}\).

Now unpack the finite \(3\)-clique condition.  If \(G\) had a clique of size
\(3\), then its three vertices could be named \(u,v,w\).  Because the clique
has cardinality \(3\), these vertices are pairwise distinct:
\[
u\ne v,\qquad u\ne w,\qquad v\ne w.
\]
Because it is a clique, each pair of distinct vertices in it is adjacent, so
\[
G.\mathrm{Adj}(u,v),\qquad G.\mathrm{Adj}(u,w),\qquad G.\mathrm{Adj}(v,w).
\]
This would be exactly the forbidden triple above.  Conversely, any pairwise
distinct \(u,v,w\) with these three adjacencies would form a \(3\)-clique on
the finite set \(\{u,v,w\}\).  Hence the no-distinct-triple assertion is
equivalent to \(G.\mathrm{CliqueFree}(3)\), and in particular it proves
\[
G.\mathrm{CliqueFree}(3).
\]

It remains to prove the independence-number bound.  The last conjunct supplied
by \(\noderef{TwoBitePositiveExistence}\), for the same \(m,p,\omega\) and the
same \(n\), is already
\[
\noderef{IndependenceNumberLess}\!\left(
(\noderef{TwoBiteFinalGraph}(\omega)).2.2,\,
(1+\varepsilon)\sqrt{(n:\mathbb R)\log(n:\mathbb R)}\right).
\]
Since \(G=(\noderef{TwoBiteFinalGraph}(\omega)).2.2\), this is precisely
\[
\noderef{IndependenceNumberLess}\!\left(
G,\,(1+\varepsilon)\sqrt{(n:\mathbb R)\log(n:\mathbb R)}\right).
\]
Thus the chosen graph \(G\) satisfies both required conclusions for this
\(n\).  Since the construction works for every \(n\ge n_0\), the theorem
follows.
\end{proof}
